% 第五章：2026-10-16 用户确认的初稿与七项扩写，待整稿按五页限制择优合并。
% 本章为素材完整缓存，不表示当前十一点全部进入最终五页稿。
\section{创新点介绍}
本方案围绕有限样条面片的解析化需求，将目标参数域中的边界建模、几何特征驱动的候选搜索和连续几何验收相结合。其技术特色主要体现在以下方面。

\paragraph{（1）以目标 NURBS 参数域中的线性 PCurve 统一描述有限边界}
本方案将线性约束明确施加于目标解析母面对应的 NURBS 参数域，通过参数直线段与母面映射的复合构造三维边界。在这一表示下，等参边界与斜向直线边界具有统一的数学描述，可以共同纳入有限参数区域的恢复过程。

这一设计将解析化对象从母面类型及其几何参数，进一步扩展至具有明确边界的有限面片。源、目标边界允许采用连续单调的参数对应，因此能够在验收中处理参数化差异，避免将相同参数值处的位置差异直接等同于几何形状差异。

\paragraph{（2）以几何特征驱动分层搜索，实现专用识别与一般优化的协同}
本方案将等参结构识别作为优先路径，将一般母面候选搜索与参数优化作为补充路径。专用分支利用明确的等参几何结构缩小求解范围；一般分支通过多类型候选生成、去重与排序，处理结构信息不足或专用验收未通过的输入。

其智能决策体现为依据几何特征、候选评价和验收反馈选择后续计算：可直接通过验收的候选无需优化，需要改善的候选再进入参数求解，仍不满足条件则继续搜索。两条路径共享总时间预算，使计算资源分配与候选质量、求解进展相联系。

\paragraph{（3）以正比例切向残差构造连续验收机制}
针对边界参数化差异引起的切向量尺度变化，本方案采用正比例缩放残差构造角度上界。设同一局部参数区间上的源、目标切向量或其有理表示的切向分子分别为 $\mathbf A$、$\mathbf B$，选取连续为正的缩放函数 $\lambda$，构造
\[\mathbf R=\mathbf A-\lambda\mathbf B.\]
当源切向长度具有正下界，且残差范数上界与该下界之比小于 $1$ 时，可据此确认切向同向并建立夹角正弦的保守上界。该构造将方向差异与参数速度差异分离，避免直接构造切向样条叉乘。

实现中结合 Bernstein 系数界、带舍入误差包络的快速判定、区间重算与自适应细分，在连续区间上进行验收。缩放函数的数值选择仅用于构造候选残差，最终接受仍依赖保守界。一般分支进一步针对实际导出的三维边界曲线验收，使判定覆盖最终几何表示。无法获得足够紧的上界时，继续细分或返回未认证状态，不将估计未通过直接解释为真实误差超限。

\paragraph{（4）以显式退化判据和可追溯失败状态增强工程鲁棒性}
本方案将极点、顶点及微小边界段作为需要专门处理的几何情形。对于微小分段，仅在源段与目标段均满足连续三维直径上界条件时允许切向豁免，同时保留完整的位置验收并记录对应区间。由此，退化处理具有明确的适用条件和覆盖范围，避免通过删除边界区间或放宽位置容差获得表面上的成功。

在执行层面，方案区分几何超差、无法认证和预算耗尽等状态，并保留分支路径与阶段诊断信息。超过单例总耗时上限的结果统一判为计算失败，使几何正确性、验收完整性与计算成本共同构成结果接受条件。

\paragraph{（5）基于齐次 Bernstein 复合的边界直接构造}
本方案利用线性 PCurve 的结构，将目标母面的二维有理表示沿参数直线进行复合，直接构造一维有理边界。设母面在局部参数单元中的齐次表示为
\[\mathbf Q(u,v)=\frac{\mathbf P(u,v)}{w(u,v)},\]
则沿线性 PCurve $\boldsymbol{\ell}(t)$ 的边界为
\[\widehat{\mathbf c}(t)=\frac{\mathbf P(\boldsymbol{\ell}(t))}{w(\boldsymbol{\ell}(t))}.\]
由于参数分量关于 $t$ 为仿射函数，分子与分母可以在 Bernstein 基下通过代数复合构造；跨越参数单元的边界则按交点分段处理。

这一设计使线性 PCurve 假设直接参与输出几何的生成，减少通过离散点重新拟合边界所引入的自由度。构造过程保留有理权重和分段关系，为后续连续误差求界提供一致的数学表示。浮点实现中的系数转换及拼接误差仍通过独立验收控制。

\paragraph{（6）面向实际内核输出对象的完整验收}
本方案将验收对象延伸至几何内核实际创建的边界曲线。目标曲线完成次数统一、权重处理、分段拼接和共享端点设置后，算法重新读取内核对象的控制顶点、节点及权重，并据此构造验收数据。

由此形成两层相互衔接的检查：一层约束实际输出曲线相对于理想母面轨迹的偏差；另一层检查源边界与实际输出边界之间的位置和切向一致性。系数舍入、端点统一及内核表示规范化等操作产生的影响，都进入最终验收范围。

这一机制将“理论上构造出的边界”与“下游实际使用的边界”连接起来，同时通过共享端点组织保证相邻输出边界的端点坐标一致，为后续拓扑装配提供几何基础。

\paragraph{（7）具有下界依据的渐进评分与候选延后求值}
一般候选搜索采用由粗到细的评分机制。算法首先在部分评价点上计算最大距离，再根据候选竞争情况补充其余内部点和边界点的评价。

该机制利用一个明确的单调关系：\textbf{评价子集上的最大距离不大于完整评价集上的最大距离。}当某候选的部分评分已严格劣于当前批次的完整评分阈值时，可以延后其剩余求值，将计算资源优先用于更有希望的候选；后续批次仍可继续处理尚未消费的候选。

实现中特别区分最大距离与均方根误差：部分点上的均方根误差不具有同样的下界性质，因此不用于上述判定。该设计使评分加速具有可解释的数学依据，同时保持候选评分与最终几何验收的职责分离。

\paragraph{（8）由位置—法向约束驱动的圆环面轴线候选恢复}
针对作为扩展目标的规则圆环面，本方案根据四组采样位置与法向信息建立轴线约束。采用方向与矩组成的直线坐标表示
\[\mathbf L=(\mathbf d,\mathbf m),\qquad\mathbf d\cdot\mathbf m=0,\]
先通过带列主元的 Householder QR 分解求取约束零空间，再利用直线坐标的二次关系求解轴线候选。

对于每条有效候选轴线，算法将采样点转换至相应的子午截面坐标，通过经过中心化和尺度归一化的圆拟合恢复圆环面中心及主、次半径，并检查规则圆环面的参数条件。

这一设计将圆环面初始化分解为轴线候选枚举与子午圆恢复两个具有几何意义的步骤。存在多个候选解时保留后续评分与验收机会，避免过早将不确定的轴线选择固化为唯一结果。

\paragraph{（9）将数值可辨识性纳入候选生成条件}
有限面片可能只覆盖母面的局部区域，其位置与法向变化不足以稳定确定全部解析参数。本方案在候选生成阶段同时检查几何可辨识性和数值条件：通过局部坐标平移与尺度归一化减小坐标量级的影响，再利用奇异值比例、方向分离度及构造残差识别秩亏或轴向歧义。

对于相关方向与法向拟合，直接对原始矩阵进行分解，避免构造法方程带来的条件数平方效应。候选生成还结合分布较广的多点支撑集与多组少量点初始化，使搜索不完全依赖某一次局部采样。

这一设计把“能否稳定确定母面参数”作为候选成立的前提，并为不适定情况提供明确诊断，防止数值不稳定的参数仅因有限采样误差较小而进入后续流程。

\paragraph{（10）参数图选择与边界对应关系的协同恢复}
同一解析母面可以具有不同的参数坐标框架；周期接缝和极点的位置也会影响边界恢复。本方案将母面几何参数与参数图选择分别处理。对于球面，在保持球心和半径不变的前提下，根据边界几何生成有限数量的坐标框架候选，并逐一尝试恢复目标参数域中的线性边界。

对周期方向，边界连接采用整周期平移保持参数展开的一致性；在连续验收过程中，细分区间还可根据源曲线的几何位置更新目标参数分割点。每次对应关系改变后，重新构造局部误差与切向残差，保证验收数据始终对应当前的参数映射。

这一机制将参数化选择纳入求解过程，降低单一坐标框架或固定参数配对对恢复结果的限制。其搜索范围由候选数量和计算预算控制，不将有限搜索解释为对所有参数化的穷尽。

\paragraph{（11）在保持验收阈值不变的前提下自适应分配认证开销}
连续边界验收采用分层计算策略：首先使用带保守舍入误差包络的双精度 Bernstein 运算进行快速判定；无法判定时，转入向外舍入的区间重算，并按需细分。两层计算针对同一源、目标区间，使用一致的验收阈值。

实现还根据实际判定效果调整计算路径。当某条边上快速层与区间层反复出现位置通过、切向界仍不充分的情况时，可关闭该边收益有限的快速过滤，避免重复执行两套计算。对于有理位置误差，则融合分子差与分母的 Bernstein 乘积计算，并复用与次数有关的组合系数。

此外，位置验收得到的误差界还可用于误差导数切向快速判据；该判据不足时，再使用正比例切向残差。这使位置与切向验收能够共享已有信息，在保持各项接受条件不变的前提下组织计算开销。

% 创新归属：以上强调本问题中的方法组合与工程实现，不将 SVD、QR 等已有工具本身表述为原创算法。
% 依据：GeneralRationalBoundaryCurves.cpp、GeneralClassification.cpp、GeneralSeedNumerics.cpp、GeneralRationalBoundaryValidation.cpp。
